\section{Discussion and Conclusion}

%%PA.5.17: TODO this section? :)

%%PA.5.17: confusing the paper ends with a related work section that's living in the Discussion/Conclusion; make this its own section somewhere?

In this paper, we proposed the problem of transfer learning for CATE estimation. One immediate question the reader may be left with is why we chose the transfer learning techniques we did. We only considered two common types of transfer: 
(1) Basic fine tuning and weights sharing techniques common in the computer vision literature \citep{finetune0,finetune2,finetune1, donahu, koch},
(2) Techniques for learning an initialization that can be quickly optimized \citep{maml,hugo,reptile}. However, many further techniques exist. Yet, transfer learning is an extensively studied and perennial problem \citep{schmid,bengio,thurn2,thrun,taylor,silver}. In \cite{oriol}, the authors attempt to combine feature embeddings that can be utilized with non-parametric methods for transfer. \cite{proto} is an extension of this work that modifies the procedure for sampling examples from the support set during training. \cite{marcin} and related techniques try to meta-learn an optimizer that can more quickly solve new tasks. \cite{progressive} attempts to overcome forgetting during transfer by systematically introducing new network layers with lateral connections to old frozen layers. \cite{yu} uses networks with memory to adapt to new tasks. We invite the reader to review \cite{maml} for an excellent overview of the current transfer learning landscape. Though the majority of the discussed techniques could be extended to CATE estimation, our implementations of \cite{progressive, marcin} proved difficult to tune and consequently learned very little. Furthermore, we were not able to successfully adapt \cite{proto} to the problem of regression. We decided to instead focus our attention on algorithms for obtaining good initializations, which were easy to adapt to our problem and quickly delivered good results without extensive tuning. 

%Though our findings were largely positive, we are still left with several open questions. Most immediately, is it possible to do transfer if the input features are dynamic? This would open up a wealth of datasets wherein the inputs are similar, but not identical. YOUR MOM IS FAT



%SRK18: * We know that there are cases when one learner e.g., the T-learner or the S-learner is better. We see in our experiment that the T-learner can benefit hugely from meta learnerning. The S-learner did already do very well and it was difficult to improve its performance.

%The transfer learning techniques we investigated all involved fairly straightforward weight sharing techniques or techniques for meta learning good initialization weights. However, many other transfer learning algorithms exist. \cite{proto,progressive,schmid,donahu}